\documentclass[10pt,a4paper]{article}
\usepackage[T1]{fontenc}
\usepackage{newtxtext,newtxmath}
\usepackage{microtype}
\usepackage{amsmath}
\usepackage{siunitx}
\usepackage[margin=18mm]{geometry}
\usepackage{booktabs,longtable,array}
\usepackage{hyperref}
\usepackage{caption}
\sisetup{per-mode=symbol,detect-all}
\setlength{\emergencystretch}{1em}
\hypersetup{hidelinks,pdftitle={Technical supplement: CMJRecoveryDynamicsBench V3},pdfauthor={Author names pending owner confirmation}}
\urlstyle{same}
\captionsetup{font=small,labelfont=bf,labelsep=period,justification=raggedright,singlelinecheck=false}

\title{Technical supplement\\[0.5em]\large Retrospective reconstruction and numerical qualification of historical CMJ recovery research}
\author{Owner-review supplement}
\date{6 November 2026}

\begin{document}
\maketitle

This supplement records the detailed formulation, architecture, route-level qualification, synthetic result, uncertainty, and analysis-contract evidence summarized in the main article. Public evidence files are copied from the v1.1.0 release-linked record with their original source hashes in the source manifest. The original neural checkpoints, fitted values, and restricted records are not distributed here. Table values below are rounded for display; the accompanying CSV files retain the full stored precision.

\section{Historical formulation identities and ancestry}

The eight formulations are distinct identities. A parent is included only when the lineage register records one. “Unresolved” is not a parent edge. Formulation status is descriptive and does not rank scientific validity.

\footnotesize
\setlength{\tabcolsep}{2.6pt}
\renewcommand{\arraystretch}{1.12}
\begin{longtable}{>{\raggedright\arraybackslash}p{0.95in}>{\raggedright\arraybackslash}p{1.10in}>{\raggedright\arraybackslash}p{1.50in}>{\raggedright\arraybackslash}p{1.10in}>{\raggedright\arraybackslash}p{1.15in}}
\caption{Complete formulation matrix from the recorded lineage and methods register. A missing direct parent is shown as unresolved, not inferred.}\label{supp:formulations}\\
\toprule
Identity & Status / parent & Response / observation & Horizon / task & Limitation \\
\midrule
\endfirsthead
\toprule
Identity & Status / parent & Response / observation & Horizon / task & Limitation \\
\midrule
\endhead
\midrule
\multicolumn{5}{r}{Continued on next page}\\
\endfoot
\bottomrule
\endlastfoot
\texttt{Initial preseason} & SUPERSEDED; parent: Unresolved & event time adaptation recovery; camp history & Force and impulse innovations at H72 / D7 & Superseded sample/surface; direct parent is unresolved; task and world law match the canonical state. \\
\texttt{Canonical preseason} & SUPERSEDED; parent: Initial preseason & event time adaptation recovery; camp history & Force and impulse innovations at H72 / D7 & Canonical query surface and sample; not a new task or world law. \\
\texttt{Rich-history camp} & RETIRED; parent: Unresolved & fitness fatigue impulse response; rich monitoring & Force and impulse innovations at H72 / D7 & Retired formulation; its direct scientific parent is unresolved. \\
\texttt{Preliminary post-exposure} & SUPERSEDED; parent: Unresolved & biexponential episode response; episode summary & Force and net-impulse innovations at H24 / H48 / H72 & Direct parent is unresolved; preliminary measurement contract without a frozen force–impulse identity. \\
\texttt{Phase-consistent} & SUPERSEDED; parent: Preliminary post-exposure & biexponential episode response; phase consistent force impulse & Force and net-impulse innovations at H24 / H48 / H72 & Observation-law repair; the episode response law is retained. \\
\texttt{Correlated exposure} & SUPERSEDED; parent: Phase-consistent & correlated exposure response; phase consistent force impulse & Force and net-impulse innovations at H24 / H48 / H72 & Material correlated exposure-distribution change; later redesigned. \\
\texttt{Threshold response} & PROPOSED; parent: Correlated exposure & threshold response; phase consistent force impulse & Force and net-impulse innovations at H24 / H48 / H72 & Proposed alternate, not an accepted successor formulation. \\
\texttt{Fixed-mode discrepancy} & ACTIVE CANDIDATE; parent: Correlated exposure & fixed mode discrepancy response; phase consistent force impulse & Force and net-impulse innovations at H24 / H48 / H72 & Active candidate; fixed modes are mathematical anchors, not biological compartments. \\
\end{longtable}


\section{Original neural architecture and input contracts}

The architecture inventory is separate from the eight scientific formulations. Each family has two repeats, r0 and r1. Parameter totals are copied from the sanitized architecture qualification table; no private weights are included.

\footnotesize
\setlength{\tabcolsep}{2.3pt}
\renewcommand{\arraystretch}{1.12}
\begin{longtable}{>{\raggedright\arraybackslash}p{0.80in}>{\raggedright\arraybackslash}p{1.05in}>{\raggedright\arraybackslash}p{3.25in}r}
\caption{Sanitized architecture and inference contract by family. Repeat-specific private weights and hashes are excluded.}\label{supp:architecture}\\
\toprule
Model / repeats & Inputs & Architecture & Parameters \\
\midrule
\endfirsthead
\toprule
Model / repeats & Inputs & Architecture & Parameters \\
\midrule
\endhead
\midrule
\multicolumn{4}{r}{Continued on next page}\\
\endfoot
\bottomrule
\endlastfoot
\texttt{M0 (r0/r1)} & 39 raw features & 39 → 32 → 2; exact GELU & 1,346 \\
\texttt{M1 (r0/r1)} & 39 raw features & 39 → 64 → 2; exact GELU & 2,690 \\
\texttt{M2 (r0/r1)} & 39 raw features & 39 → 64 → 64 → 2; exact GELU & 6,850 \\
\texttt{M3 (r0/r1)} & 39 raw + 4 × 4 history & Raw backbone + pointwise history encoder; mean/latest pooling → 128 → 2 & 7,138 \\
\texttt{M4 (r0/r1)} & 39 raw + 4 × 4 history & Raw backbone + 4 → 32 GRU; mean/latest pooling → 128 → 2 & 10,626 \\
\texttt{M5 (r0/r1)} & 39 raw + 4 × 4 history + horizon & M4 + 2 × 8 horizon embedding → 136 → 2 & 10,658 \\
\texttt{M6 (r0/r1)} & 39 raw + 4 × 4 history + horizon + 5 basis values & M5 + first two basis values → 16 → 152 → 2 & 10,738 \\
\texttt{M7 (r0/r1)} & 39 raw + 4 × 4 history + horizon + 5 basis values & M5 + all five basis values → 16 → 152 → 2 & 10,786 \\
\texttt{M8 (r0/r1)} & 39 raw + 4 × 4 history + horizon + 5 basis values & M7 + bias-free 5 → 2 linear basis head & 10,796 \\
\end{longtable}


\section{Route-by-route numerical qualification}

The stored saved-CUDA rule is \(|a_i-b_i|\leq 10^{-5}+10^{-4}|b_i|\) for every compared element. The CPU-source comparison separately requires exact array equality. The reported maximum absolute discrepancy combines two outputs with different physical units and is not a comparable physical effect size. “Not tabulated” is retained when the source record did not report a failing-element count for a passing route.

\small
\setlength{\tabcolsep}{4pt}
\renewcommand{\arraystretch}{1.12}
\begin{longtable}{>{\raggedright\arraybackslash}p{0.70in}>{\raggedright\arraybackslash}p{0.85in}>{\raggedright\arraybackslash}p{0.85in}p{0.95in}p{1.40in}p{1.35in}}
\caption{All 18 sanitized model-repeat records. The recorded inference-rerun flag is false for every route. The maximum absolute error pools force and impulse outputs with different units and is descriptive only.}\label{supp:routes}\\
\toprule
Route & CPU-source & Saved CUDA & ATOL/RTOL & Max abs. error & Failing elements \\
\midrule
\endfirsthead
\toprule
Route & CPU-source & Saved CUDA & ATOL/RTOL & Max abs. error & Failing elements \\
\midrule
\endhead
\midrule
\multicolumn{6}{r}{Continued on next page}\\
\endfoot
\bottomrule
\endlastfoot
\texttt{M0-r0} & Exact & Pass & 1e-05 / 1e-04 & 1.43051147e-06 & 0 \\
\texttt{M0-r1} & Exact & Pass & 1e-05 / 1e-04 & 1.43051147e-06 & 0 \\
\texttt{M1-r0} & Exact & Pass & 1e-05 / 1e-04 & 1.43051147e-06 & not tabulated \\
\texttt{M1-r1} & Exact & Pass & 1e-05 / 1e-04 & 1.43051147e-06 & not tabulated \\
\texttt{M2-r0} & Exact & Pass & 1e-05 / 1e-04 & 1.90734863e-06 & not tabulated \\
\texttt{M2-r1} & Exact & Pass & 1e-05 / 1e-04 & 1.90734863e-06 & not tabulated \\
\texttt{M3-r0} & Exact & Pass & 1e-05 / 1e-04 & 2.38418579e-06 & not tabulated \\
\texttt{M3-r1} & Exact & Pass & 1e-05 / 1e-04 & 1.19209290e-06 & not tabulated \\
\texttt{M4-r0} & Exact & FAIL & 1e-05 / 1e-04 & 4.75883484e-04 & 254 \\
\texttt{M4-r1} & Exact & FAIL & 1e-05 / 1e-04 & 3.26395035e-04 & 146 \\
\texttt{M5-r0} & Exact & FAIL & 1e-05 / 1e-04 & 3.71456146e-04 & 443 \\
\texttt{M5-r1} & Exact & FAIL & 1e-05 / 1e-04 & 2.70366669e-04 & 206 \\
\texttt{M6-r0} & Exact & FAIL & 1e-05 / 1e-04 & 1.99317932e-04 & 88 \\
\texttt{M6-r1} & Exact & FAIL & 1e-05 / 1e-04 & 2.79307365e-04 & 140 \\
\texttt{M7-r0} & Exact & FAIL & 1e-05 / 1e-04 & 3.17573547e-04 & 163 \\
\texttt{M7-r1} & Exact & FAIL & 1e-05 / 1e-04 & 3.04937363e-04 & 198 \\
\texttt{M8-r0} & Exact & FAIL & 1e-05 / 1e-04 & 2.21252441e-04 & 98 \\
\texttt{M8-r1} & Exact & FAIL & 1e-05 / 1e-04 & 3.04222107e-04 & 145 \\
\end{longtable}


All 18 routes have recorded exact CPU-source equality. The eight M0--M3 saved-CUDA comparisons pass and the ten M4--M8 comparisons fail. These are preserved qualification results, not a fresh model inference or GPU replay.

\section{Full public synthetic six-cell results}

Each row shows a stored cellwise NRMSE, which divides RMSE by that cell's population standard deviation (denominator convention \(n_c\), not \(n_c-1\)). Root A, B, and C are three caller-supplied public evaluation roots from the same generator. The mean and sample standard deviation summarize those three roots; they are not a confidence interval or human-population estimate.

\footnotesize
\setlength{\tabcolsep}{2.0pt}
\renewcommand{\arraystretch}{1.12}
\begin{longtable}{>{\raggedright\arraybackslash}p{0.80in}>{\raggedright\arraybackslash}p{1.62in}rrrrr}
\caption{Cellwise NRMSE by public evaluation root. Values are rounded to 6 decimal places for display; full stored precision is available in the accompanying numerical data.}\label{supp:sixcell}\\
\toprule
Model & Outcome-horizon cell & Root A & Root B & Root C & Mean & Sample SD \\
\midrule
\endfirsthead
\toprule
Model & Outcome-horizon cell & Root A & Root B & Root C & Mean & Sample SD \\
\midrule
\endhead
\midrule
\multicolumn{7}{r}{Continued on next page}\\
\endfoot
\bottomrule
\endlastfoot
Zero innovation & H24 / net impulse (m/s) & 2.757823 & 2.759502 & 2.767942 & 2.761756 & 0.005423 \\
Training mean & H24 / net impulse (m/s) & 1.053386 & 1.036381 & 1.045551 & 1.045106 & 0.008512 \\
Ridge regression & H24 / net impulse (m/s) & 0.651329 & 0.651082 & 0.654060 & 0.652157 & 0.001653 \\
Structured model & H24 / net impulse (m/s) & 0.694226 & 0.698253 & 0.695487 & 0.695989 & 0.002060 \\
Zero innovation & H24 / force innovation (N/kg) & 3.043088 & 3.049200 & 3.008360 & 3.033549 & 0.022028 \\
Training mean & H24 / force innovation (N/kg) & 1.054776 & 1.040758 & 1.047428 & 1.047654 & 0.007012 \\
Ridge regression & H24 / force innovation (N/kg) & 0.630319 & 0.630139 & 0.612765 & 0.624408 & 0.010083 \\
Structured model & H24 / force innovation (N/kg) & 0.663259 & 0.664902 & 0.644444 & 0.657535 & 0.011367 \\
Zero innovation & H48 / net impulse (m/s) & 2.292457 & 2.222562 & 2.278043 & 2.264354 & 0.036904 \\
Training mean & H48 / net impulse (m/s) & 1.035905 & 1.022505 & 1.036852 & 1.031754 & 0.008024 \\
Ridge regression & H48 / net impulse (m/s) & 0.752218 & 0.765857 & 0.761406 & 0.759827 & 0.006955 \\
Structured model & H48 / net impulse (m/s) & 0.797196 & 0.812503 & 0.804108 & 0.804603 & 0.007666 \\
Zero innovation & H48 / force innovation (N/kg) & 2.525989 & 2.550929 & 2.549185 & 2.542034 & 0.013923 \\
Training mean & H48 / force innovation (N/kg) & 1.044715 & 1.028447 & 1.045208 & 1.039457 & 0.009538 \\
Ridge regression & H48 / force innovation (N/kg) & 0.737124 & 0.737851 & 0.720804 & 0.731926 & 0.009639 \\
Structured model & H48 / force innovation (N/kg) & 0.767842 & 0.768971 & 0.746271 & 0.761028 & 0.012792 \\
Zero innovation & H72 / net impulse (m/s) & 1.882918 & 1.844269 & 1.892503 & 1.873230 & 0.025535 \\
Training mean & H72 / net impulse (m/s) & 1.028582 & 1.018947 & 1.026964 & 1.024831 & 0.005159 \\
Ridge regression & H72 / net impulse (m/s) & 0.834328 & 0.829607 & 0.836870 & 0.833602 & 0.003686 \\
Structured model & H72 / net impulse (m/s) & 0.877276 & 0.876384 & 0.879950 & 0.877870 & 0.001856 \\
Zero innovation & H72 / force innovation (N/kg) & 2.110591 & 2.101190 & 2.115307 & 2.109029 & 0.007187 \\
Training mean & H72 / force innovation (N/kg) & 1.032853 & 1.019695 & 1.029735 & 1.027428 & 0.006875 \\
Ridge regression & H72 / force innovation (N/kg) & 0.814940 & 0.821730 & 0.810435 & 0.815702 & 0.005686 \\
Structured model & H72 / force innovation (N/kg) & 0.843228 & 0.850143 & 0.835237 & 0.842869 & 0.007459 \\
\end{longtable}


\section{Recorded camp-bootstrap intervals}

The stored percentile intervals use 5,000 paired bootstrap replicates per root. The resampling unit is camp: 16 camps are sampled with replacement and the 250 episodes and three horizons for each selected camp remain together. The paired contrast is model minus zero innovation; negative values favor the named model over that baseline. The receipt did not record a ridge-versus-structured interval.

\footnotesize
\setlength{\tabcolsep}{2.1pt}
\renewcommand{\arraystretch}{1.12}
\begin{longtable}{p{0.42in}p{1.15in}rrrrr}
\caption{Recorded 95\% camp-bootstrap percentile endpoints by root. Delta is model minus zero innovation. No ridge-versus-structured interval is recorded.}\label{supp:uncertainty}\\
\toprule
Root & Model & NRMSE low & NRMSE high & Delta low & Delta high & Replicates \\
\midrule
\endfirsthead
\toprule
Root & Model & NRMSE low & NRMSE high & Delta low & Delta high & Replicates \\
\midrule
\endhead
\midrule
\multicolumn{7}{r}{Continued on next page}\\
\endfoot
\bottomrule
\endlastfoot
A & \texttt{ridge\_full} & 0.728033 & 0.746188 & -1.752167 & -1.644956 & 5000 \\
A & \texttt{structured\_full} & 0.763766 & 0.784347 & -1.714653 & -1.608125 & 5000 \\
B & \texttt{ridge\_full} & 0.733825 & 0.745350 & -1.715574 & -1.648424 & 5000 \\
B & \texttt{structured\_full} & 0.773494 & 0.784623 & -1.676622 & -1.609237 & 5000 \\
C & \texttt{ridge\_full} & 0.726382 & 0.740127 & -1.733792 & -1.669113 & 5000 \\
C & \texttt{structured\_full} & 0.761200 & 0.775063 & -1.700474 & -1.632528 & 5000 \\
\end{longtable}


\section{Ablation and nested-camp result detail}

The ablation variants are distinct receipt-locked model configurations. The structured-without-history variant is a separate comparator diagnostic. The nested-camp analysis uses 24, 48, and 96 training camps from a single recorded seeded ordering; the alpha selected on full training/development data is held fixed. Root values are descriptive and do not identify causal feature effects.

\footnotesize
\setlength{\tabcolsep}{2.0pt}
\renewcommand{\arraystretch}{1.12}
\begin{longtable}{>{\raggedright\arraybackslash}p{1.30in}r>{\raggedright\arraybackslash}p{2.45in}>{\raggedright\arraybackslash}p{1.45in}}
\caption{Recorded six-cell NRMSE root values for selected ablation and nested-camp comparisons. Values are rounded to 6 decimal places; full precision is retained in the accompanying numerical data.}\label{supp:sample-efficiency}\\
\toprule
Variant & Training camps & Evaluation-root NRMSE & Root mean ± sample SD \\
\midrule
\endfirsthead
\toprule
Variant & Training camps & Evaluation-root NRMSE & Root mean ± sample SD \\
\midrule
\endhead
\midrule
\multicolumn{4}{r}{Continued on next page}\\
\endfoot
\bottomrule
\endlastfoot
Full ridge & 96 & A 0.736710; B 0.739378; C 0.732724 & $0.736270\pm0.003349$ \\
Ridge without prior history & 96 & A 0.822457; B 0.823679; C 0.812906 & $0.819681\pm0.005899$ \\
Ridge without current exposure & 96 & A 0.868359; B 0.861227; C 0.873620 & $0.867735\pm0.006220$ \\
Ridge without participant context & 96 & A 0.736746; B 0.739455; C 0.732837 & $0.736346\pm0.003327$ \\
Structured without prior history & 96 & A 0.823035; B 0.824241; C 0.813749 & $0.820342\pm0.005741$ \\
Ridge, nested 24 camps & 24 & A 0.736637; B 0.740818; C 0.733212 & $0.736889\pm0.003809$ \\
Ridge, nested 48 camps & 48 & A 0.736861; B 0.739726; C 0.733219 & $0.736602\pm0.003261$ \\
Ridge, nested 96 camps & 96 & A 0.736710; B 0.739378; C 0.732724 & $0.736270\pm0.003349$ \\
Structured, nested 24 camps & 24 & A 0.774179; B 0.778755; C 0.767939 & $0.773624\pm0.005429$ \\
Structured, nested 48 camps & 48 & A 0.774324; B 0.778711; C 0.767828 & $0.773621\pm0.005475$ \\
Structured, nested 96 camps & 96 & A 0.773838; B 0.778526; C 0.767583 & $0.773316\pm0.005490$ \\
\end{longtable}


\section{Public synthetic study protocol}

\begin{longtable}{>{\raggedright\arraybackslash}p{1.45in}>{\raggedright\arraybackslash}p{5.15in}}
\toprule
Protocol item & Recorded value \\
\midrule
\endfirsthead
\toprule
Protocol item & Recorded value \\
\midrule
\endhead
Training split & 96 camps; 24,000 episodes; 72,000 horizon rows. \\
Development split & 16 camps; 4,000 episodes; 12,000 horizon rows. \\
Each evaluation root & 16 camps; 4,000 episodes; 12,000 horizon rows; 4,000 rows per outcome-horizon cell. \\
Generator relation & Three caller-supplied public roots from one fixed-mode synthetic generator; no human data or private historical test root. \\
Predictor schema & 82 public source fields; one-hot horizon expansion to 84 numeric features; grouping keys, target labels, target-time realizations, hidden effects, and oracle states excluded. \\
Preprocessing & Training-only means and population standard deviations; zero scales replaced by one. \\
Models & Zero innovation; training-only outcome-by-horizon mean; cellwise ridge; structured comparator using observed predictors and fixed 24 h/84 h mathematical modes. \\
Alpha selection & Grid \(\{0.01,0.1,1,10,100\}\); minimum development six-cell mean NRMSE; lower alpha breaks ties; all full-data selections choose 0.01. \\
Registered cells & H24/H48/H72 crossed with the registered force-innovation cell and net-impulse innovation. \\
Primary aggregate & Equal-weight mean of six cellwise NRMSE values; each cell is normalized separately. \\
Bootstrap & PCG64; 5,000 paired whole-camp resamples; 16 camps per root; seed 49,505,000; all episode and horizon records stay with the camp. \\
Sample-size analysis & 24, 48, 96 camps from one seeded nested permutation; alpha chosen on the full dataset held fixed; seed 52,610,526. \\
\bottomrule
\end{longtable}

\section{Native measurement and evaluation contract}

For the native task, force is reported in \(\mathrm{N\,kg^{-1}}\) and gravity-adjusted net impulse in \(\mathrm{m\,s^{-1}}\). With a supplied force trace \(F_z(t)\), system mass \(m\), gravity \(g\), concentric onset \(t_{\mathrm{co}}\), takeoff \(t_{\mathrm{to}}\), and duration \(T=t_{\mathrm{to}}-t_{\mathrm{co}}\):
\begin{align}
Y_F &= \frac{1}{mT}\int_{t_{\mathrm{co}}}^{t_{\mathrm{to}}}F_z(t)\,\mathrm{d}t,\\
Y_J &= \frac{1}{m}\int_{t_{\mathrm{co}}}^{t_{\mathrm{to}}}\left(F_z(t)-mg\right)\,\mathrm{d}t.
\end{align}
The implementation interpolates force at event boundaries and uses composite trapezoids. It does not detect events or authenticate calibration, mass, or trial validity. An assessment is the arithmetic mean of exactly three distinct valid trials. Baseline is the mean of the two latest distinct valid assessments in the inclusive 24--240 h pre-origin interval. Native target windows are H72 at 66--78 h and D7 at 156--180 h.

The four native analysis cells are force and impulse at H72 and D7. Within each cell, query-weighted RMSE, mean absolute error (MAE), and signed bias are computed separately. Repeats are reported separately; a family metric is the arithmetic mean of repeat metrics, never an ensemble of predictions. Uncertainty resamples participant groups, preserving all participant queries, origins, horizons, candidates, repeats, and baselines. The 95\% percentile intervals are marginal.

The native primary sign-flip test uses participant sums of zero-baseline squared error minus mean two-repeat squared error. The one-sided null requires independent participant groups and symmetry of cluster loss differences. At \(n\leq16\), all \(2^n\) sign assignments are enumerated; above that threshold, sampled assignments use a plus-one Monte Carlo correction. Holm adjustment applies to four locked primary tests. No athlete-level result from this analysis was generated for this manuscript. The stored six-participant engineering rehearsal uses invented values only; its smallest exact \(p\) value is \(1/64\), and the four-cell Holm-adjusted value is \(4/64=0.0625\), so it cannot pass a 0.05 family-wise threshold.

\section{Evidence and interpretation boundaries}

\begin{itemize}
\item The public scientific software identity is release v1.1.0 at commit \texttt{401f87aad1e4ef4108011fa0c1d769a24f1ad329}.
\item The M8 synthetic-study result and reproduction receipts keep their own package-source, protocol, selection-lock, and result digests. They are not relabeled as a new run of the original neural model campaign.
\item Private checkpoint bytes, fitted preprocessing values, and restricted records are excluded. The public-safe route summary supports figure regeneration but not independent replay of those artifacts.
\item No new training, fit, checkpoint inference, empirical scoring, athlete-data acquisition, or GPU replay was performed for this manuscript.
\item Historical saved-CUDA failures for M4--M8 and rich-history row-count discrepancies remain negative findings. No cause for the CUDA discrepancy is asserted.
\item No independent athlete predictive validity, clinical utility, out-of-distribution transportability, or production readiness is established.
\end{itemize}

\end{document}
